\documentclass[11pt,reqno]{amsart}

\usepackage[T1]{fontenc}
\usepackage{lmodern}
\usepackage{microtype}
\usepackage{geometry}
\usepackage{amsmath,amssymb,mathtools}
\usepackage{booktabs,longtable,array}
\usepackage{enumitem}
\usepackage{xcolor}
\usepackage{hyperref}
\usepackage[nameinlink,capitalize,noabbrev]{cleveref}

\geometry{margin=1.05in}
\hypersetup{colorlinks=true,linkcolor=blue!55!black,citecolor=green!35!black,
  urlcolor=blue!65!black,
  pdftitle={Formalization-ready proof packet for dyadic Galois presentations},
  pdfauthor={OpenAI, prepared for David Roe}}
\setlist{itemsep=0.25em,topsep=0.45em}
\allowdisplaybreaks

\newtheorem{theorem}{Theorem}[section]
\newtheorem{proposition}[theorem]{Proposition}
\newtheorem{lemma}[theorem]{Lemma}
\newtheorem{corollary}[theorem]{Corollary}
\newtheorem{definition}[theorem]{Definition}
\newtheorem{construction}[theorem]{Construction}
\theoremstyle{remark}
\newtheorem{remark}[theorem]{Remark}
\newtheorem{warning}[theorem]{Warning}

\numberwithin{equation}{section}

\newcommand{\Q}{\mathbf Q}
\newcommand{\Z}{\mathbf Z}
\newcommand{\F}{\mathbf F}
\newcommand{\ur}{\mathrm{ur}}
\newcommand{\tame}{\mathrm{t}}
\newcommand{\pc}{\mathrm{pc}}
\newcommand{\sgn}{\mathrm{sgn}}
\newcommand{\prof}{\mathrm{prof}}
\newcommand{\Gal}{\operatorname{Gal}}
\newcommand{\Hom}{\operatorname{Hom}}
\newcommand{\Sur}{\operatorname{Sur}}
\newcommand{\Aut}{\operatorname{Aut}}
\newcommand{\im}{\operatorname{im}}
\newcommand{\coker}{\operatorname{coker}}
\newcommand{\inv}{\operatorname{inv}}
\newcommand{\Tr}{\operatorname{Tr}}
\newcommand{\Frat}{\operatorname{Frat}}
\newcommand{\whZ}{\widehat{\Z}}
\newcommand{\Gsum}{\mathcal G}
\newcommand{\cU}{\mathcal U}
\newcommand{\cC}{\mathcal C}
\newcommand{\cS}{\mathcal S}
\newcommand{\cW}{\mathcal W}
\newcommand{\eps}{\varepsilon}
\DeclareMathOperator{\rad}{rad}
\DeclareMathOperator{\Otwo}{O_2}

\title[Formalization packet for dyadic presentations]
{A Formalization-Ready Proof Packet for\\
Presentations of Dyadic Absolute Galois Groups}
\author{Prepared for David Roe}
\date{July 28, 2026}

\begin{document}
\begin{abstract}
This supplement audits and restructures the July 2026 draft on explicit
presentations of absolute Galois groups of dyadic local fields.  It gives
complete proofs of the boundary, exact-lifting, composition-series,
Fourier--Gauss, and local determinant reductions; corrects the arithmetic
branch decomposition; and states the remaining word and marked-Demu\v{s}kin
obligations as finite, kernel-checkable certificates.  The generic
finite-target induction is not reproved from scratch: it is already present
in the accompanying \texttt{gq2-lean} development as a source-generic theorem,
and this packet specifies the degree-\(n\) interface needed to reuse it.

The resulting theorem is unconditional once two explicit certificate
families are checked: the marked automorphism-lifting certificate for the
standard \(q=2\) Demu\v{s}kin cores, and the universal Fox--Stokes--Hessian
certificate for each displayed relator.  These are exactly the two places
where the uploaded draft currently asserts calculations that it does not
show.  No finite testing is promoted to a proof.
\end{abstract}

\maketitle
\tableofcontents

\section{Logical status and principal corrections}

Let \(K/\Q_2\) be finite, let \(n=[K:\Q_2]\), let \(q_K=2^f\), and let
\(D_K=G_K(2)\).  For a proposed wild word \(R\), write
\[
 \Gamma_R=\left\langle\sigma,\tau,x_0,\ldots,x_n\ \middle|\
 \tau^\sigma=\tau^{q_K},\ R=1,\
 \overline{\langle\!\langle x_0,\ldots,x_n\rangle\!\rangle}
 \text{ is pro-}2\right\rangle_{\prof}.
\]
The compact words must use
\[
 x_2^{-\sigma}(x_2\tau)^{\omega_2},
\]
not \(x_2^{-\sigma_2}(x_2\tau)^{\omega_2}\).  The uploaded revision has
made this correction.

There are three further structural corrections.

\begin{enumerate}
\item The even-\(\eta\) ``sign-Frobenius'' row is incompatible with the
standing hypothesis that \(K(i)/K\) is ramified; see
\cref{prop:sign-excluded}.  It should be removed from the new ramified-\(i\)
part of the theorem and, if retained at all, placed with the
unramified-\(i\) discussion.
\item Equality of one unshifted Gauss sum is not by itself the Fourier
interface needed by the exact-image induction.  The correct interface is
an affine Gauss/phase-cover interface; see \cref{sec:fourier}.
\item The two-character Labute statement in the draft is much stronger
than Labute's classification and its displayed proof is only an outline.
For formalization it should be replaced by the group-level marked matching
reduction in \cref{sec:marked}; the only new theorem there is an explicit
per-core automorphism-lifting certificate.
\end{enumerate}

\begin{theorem}[Kernel-checkable main theorem]\label{thm:certificate-main}
Fix a ramified-\(i\) field \(K\), a standard marked core \(P_K\), and a
wild word \(R_K\).  Suppose:
\begin{enumerate}
\item \(P_K\) has a marked-core certificate in the sense of
\cref{def:core-certificate};
\item \(R_K\) has a word certificate in the sense of
\cref{def:word-certificate}; and
\item the standard local inputs listed in \cref{sec:local-inputs} hold for
\(K\).
\end{enumerate}
Then \(\Gamma_{R_K}\cong G_K\).  The isomorphism identifies the tame
quotient, the maximal pro-\(2\) quotient, the cyclotomic orientation, and
the full geometric unramified character.
\end{theorem}

\begin{proof}
The marked-core certificate gives the fully marked maximal pro-\(2\)
identification.  The tame and joint boundary statements follow from
\cref{thm:boundary}.  The word certificate supplies the source-side
lifting, Stokes, scalar, and affine determinant obligations.  The local
inputs and \cref{thm:local-gauss} supply the corresponding field-side
obligations.  The source-generic exact-image theorem then gives equality of
all finite epimorphism counts.  Since both groups are topologically finitely
generated, profinite reconstruction gives the asserted isomorphism.  Each
step is expanded below.
\end{proof}

\section{Conventions for profinite words}

\begin{definition}[Profinite exponentiation]
For an element \(g\) of a profinite group, the canonical continuous map
\(\Z\to\overline{\langle g\rangle}\), \(a\mapsto g^a\), extends uniquely to
\(\whZ\).  Its value at \(a\in\whZ\) is denoted \(g^a\).  If \(x,g\) are
group elements, set
\[
 x^g=g^{-1}xg,\qquad [x,y]=x^{-1}y^{-1}xy,
 \qquad x^{-g}:=(x^{-1})^g.
\]
Let \(\omega_2\in\whZ\) have component \(1\) in \(\Z_2\) and component
\(0\) in every \(\Z_\ell\), \(\ell\ne2\), and put
\(\sigma_2=\sigma^{\omega_2}\).
\end{definition}

\begin{lemma}[Finite evaluation of \(\omega_2\)]\label{lem:omega-finite}
If \(g\) has finite order \(2^a m\) with \(m\) odd, then \(g^{\omega_2}=g^e\),
where \(e\) is the unique residue modulo \(2^a m\) satisfying
\[
 e\equiv1\pmod{2^a},\qquad e\equiv0\pmod m.
\]
In particular, \(g^{\omega_2}\) is the \(2\)-primary component of \(g\).
\end{lemma}

\begin{proof}
The assertion is the Chinese remainder description of the image of
\(\omega_2\) in \(\Z/(2^am)\Z\).  Since every continuous map from a
profinite word group to a finite group factors through a finite quotient,
this also gives a terminating evaluator for every occurrence of
\(\omega_2\).
\end{proof}

\begin{remark}[Lean encoding]
A word syntax should distinguish integral powers, \(\Z_2\)-powers, and
\(\whZ\)-powers.  Finite evaluation should reduce the last kind by
\cref{lem:omega-finite}; this avoids selecting global integer
representatives.  The expression \(x^{-g}\) should be a defined syntax
constructor or expanded to \((x^{-1})^g\), never parsed as an exponent
\(-g\).
\end{remark}

\section{The tame quotient and the common boundary}

Put
\[
 T_q=\langle\sigma,\tau\mid\tau^\sigma=\tau^q\rangle_{\prof},
 \qquad \nu_2(\sigma)=1,\quad\nu_2(\tau)=0,
\]
where \(\nu_2:T_q\to\Z_2\) is the \(2\)-primary component of the
unramified character.

\begin{lemma}[Tame inertia is pro-odd]\label{lem:tame-odd}
In every finite quotient of \(T_q\), the order of the image of \(\tau\) is
prime to \(q\).  Consequently the closed subgroup generated by all
conjugates of \(\tau\) is pro-prime-to-\(2\) when \(q\) is a power of
\(2\).
\end{lemma}

\begin{proof}
Let the image of \(\tau\) have order \(m\).  The elements \(\tau\) and
\(\tau^q\) are conjugate, hence have the same order.  But
\(\operatorname{ord}(\tau^q)=m/\gcd(m,q)\), so \(\gcd(m,q)=1\).  Passing
over all finite quotients gives the pro-prime-to-\(2\) assertion.
\end{proof}

\begin{lemma}\label{lem:ker-nu-odd}
The kernel of \(\nu_2:T_q\to\Z_2\) is pro-odd.
\end{lemma}

\begin{proof}
In a finite quotient, the image of the kernel is an extension of a finite
quotient of tame inertia, which has odd order by \cref{lem:tame-odd}, by a
finite quotient of the prime-to-\(2\) part of the unramified procyclic
group.  Both kernel and quotient have odd order.
\end{proof}

\begin{lemma}[No normal pro-\(2\) subgroup in the tame quotient]
\label{lem:O2-tame}
\(\Otwo(T_q)=1\).
\end{lemma}

\begin{proof}
Let \(N\triangleleft T_q\) be pro-\(2\).  By \cref{lem:ker-nu-odd},
\(N\cap\ker\nu_2=1\), so \(\nu_2|_N\) is injective.  If \(z\in N\), then
for every tame-inertia element \(t\), the commutator \([z,t]\) lies both
in \(N\) and in \(\ker\nu_2\), hence is trivial.  Thus \(z\) centralizes
\(\tau\).

Write \(a=\nu_2(z)\in\Z_2\).  For every \(k\ge1\), map \(T_q\) to
\[
 C_{q^{2^k}-1}\rtimes C_{2^k},
\]
where the generator of \(C_{2^k}\) acts on
\(C_{q^{2^k}-1}\) by the \(q\)-power map.  This is well-defined, because
the order of \(q\) modulo \(q^{2^k}-1\) is exactly \(2^k\): if
\(0<d<2^k\), then \(0<q^d-1<q^{2^k}-1\).  Since the image of \(z\)
centralizes the inertia generator, \(a\equiv0\pmod{2^k}\).  This holds for
all \(k\), so \(a=0\).  Injectivity of \(\nu_2|_N\) gives \(z=1\), hence
\(N=1\).
\end{proof}

Let \(W_R\) be the closed normal subgroup of \(\Gamma_R\) generated by
the \(x_i\).

\begin{proposition}[Tame and maximal pro-\(2\) specializations]
\label{prop:specializations}
Assume that killing all \(x_i\) makes the wild word \(R\) trivial, and
that killing \(\tau\) and replacing every profinite exponent by its
\(2\)-component changes \(R\) to a pro-\(2\) word \(P\).  Then:
\begin{enumerate}
\item \(W_R=\Otwo(\Gamma_R)\) and \(\Gamma_R/W_R\cong T_{q_K}\);
\item the maximal pro-\(2\) quotient of \(\Gamma_R\) is
\[
 D_P=\langle\sigma,x_0,\ldots,x_n\mid P=1\rangle_{\mathrm{pro}-2};
\]
\item the induced unramified character has
\(\nu(\sigma)=1\) and \(\nu(x_i)=0\).
\end{enumerate}
\end{proposition}

\begin{proof}
The admissibility condition makes \(W_R\) pro-\(2\), and the first
specialization identifies the quotient with \(T_{q_K}\).  Any normal
pro-\(2\) subgroup of \(\Gamma_R\) maps to a normal pro-\(2\) subgroup of
\(T_{q_K}\), which is trivial by \cref{lem:O2-tame}; therefore it lies in
\(W_R\).  This proves (1).

Every map to a pro-\(2\) group kills \(\tau\), because tame inertia is
pro-odd, and kills every odd-primary component of \(\sigma\).  Hence it
factors through the displayed specialization.  Conversely that
specialization is a pro-\(2\) quotient, proving (2).  The generator values
in (3) are immediate.
\end{proof}

\begin{theorem}[Joint boundary surjectivity]\label{thm:boundary}
Suppose \(D_P\cong D_K\) through an isomorphism preserving the full
unramified character.  Then the natural map
\[
 \Gamma_R\longrightarrow
 \partial_K:=T_{q_K}\times_{\Z_2}D_K
\]
is surjective.  The analogous map from \(G_K\) is surjective as well.
\end{theorem}

\begin{proof}
The image projects onto both factors.  Apply the relative form of
Goursat's lemma over \(\Z_2\).  Any failure of surjectivity would give a
nontrivial common quotient of
\(\ker(T_{q_K}\to\Z_2)\) and \(\ker(D_K\to\Z_2)\).  The former is pro-odd
by \cref{lem:ker-nu-odd}; the latter is pro-\(2\), because \(D_K\) is
pro-\(2\).  Their only common profinite quotient is trivial.  Hence the
image is the full fibre product.  The field case is identical, using the
wild inertia and maximal pro-\(2\) quotient of \(G_K\).
\end{proof}

\section{Exact elementary lifting and the Fox complex}

Let \(C\) be finite, let \(A\) be a finite elementary abelian normal
subgroup in an extension
\[
 1\longrightarrow A\longrightarrow B\stackrel{\pi}{\longrightarrow}C
 \longrightarrow1,
\]
and let \(\rho:\Gamma_R\to C\) be a continuous homomorphism.  The
conjugation action through \(\rho\) makes \(A\) an \(\F_2[C]\)-module.

\begin{proposition}[Literal defect formula]\label{prop:defect}
Choose lifts in \(B\) of the \(n+3\) marked generators of \(\Gamma_R\).
Evaluation of the tame and wild relators gives a defect
\(e\in A^2\).  If the generator lifts are changed by
\(a\in A^{n+3}\), then
\[
 e\longmapsto e+d^1_{R,\rho}(a),
\]
where \(d^1_{R,\rho}\) is the evaluated Fox Jacobian of the two relators.
Consequently:
\begin{enumerate}
\item a lift \(\widetilde\rho:\Gamma_R\to B\) exists if and only if the
defect class vanishes in \(\coker d^1_{R,\rho}\);
\item when nonempty, the set of raw generator lifts satisfying both
relations is a torsor under \(\ker d^1_{R,\rho}\);
\item quotienting by simultaneous conjugation by \(A\) gives the complex
\[
 0\longrightarrow A\stackrel{d^0}{\longrightarrow}A^{n+3}
 \stackrel{d^1_{R,\rho}}{\longrightarrow}A^2\longrightarrow0.
\]
\end{enumerate}
\end{proposition}

\begin{proof}
For a finite word, the assertion follows by induction from
\[
 D(uv)=D(u)+\bar uD(v),\qquad
 D(u^{-1})=-\bar u^{-1}D(u),
\]
where \(\bar u\) is the lower value in \(C\).  A profinite word is
evaluated in the finite group \(B\), so its profinite powers may be
replaced by the finite exponents of \cref{lem:omega-finite}; the same
induction applies.  The first two conclusions are the elementary facts
that an affine equation \(d^1(a)=-e\) is solvable exactly when
\([e]=0\) in the cokernel, and that its solution set is a translate of
the kernel.  Simultaneous conjugation changes the generator tuple by the
usual coboundary \(d^0\), giving the complex.
\end{proof}

\begin{proposition}[No extra linear equation from admissibility]
\label{prop:no-extra}
In the situation of \cref{prop:defect}, assume the lower images of the
\(x_i\) lie in \(\Otwo(C)\).  Every lift of an \(x_i\) lies in a finite
normal \(2\)-subgroup of \(B\).  Thus the pro-\(2\) admissibility condition
adds no equation to the elementary lifting problem.
\end{proposition}

\begin{proof}
The inverse image \(\pi^{-1}(\Otwo(C))\) is an extension of a finite
\(2\)-group by the elementary abelian \(2\)-group \(A\), hence is a finite
\(2\)-group.  It is normal in \(B\) and contains every lift of every
\(x_i\).
\end{proof}

\begin{lemma}[Exactness in coefficients]\label{lem:coeff-exact}
A short exact sequence of finite \(\F_2[C]\)-modules induces a termwise
short exact sequence of the corresponding Fox complexes.  The defect
construction, boundary restriction, and connecting maps are natural in
that sequence.
\end{lemma}

\begin{proof}
The three terms are finite direct powers of the coefficient module, and
the differentials are finite sums of action operators.  Exactness and
naturality therefore follow coordinatewise.
\end{proof}

\section{Extending Stokes duality from simple to nonsplit modules}

The draft correctly aims for a natural chain-level Stokes map, rather than
cardinality identities.  The following lemma is the precise induction it
needs.

\begin{lemma}[Composition-series extension]\label{lem:composition}
For every finite \(\F_2[C]\)-module \(A\), let
\[
 \eta_A:\cC^\bullet(A)\longrightarrow
 \Hom_{\F_2}(\cC^\bullet(A^\vee),\F_2)[-2]
\]
be a natural chain map.  Assume:
\begin{enumerate}
\item short exact coefficient sequences give commuting short exact
sequences of the source complexes and of the dual target complexes; and
\item \(\eta_V\) is a quasi-isomorphism for every simple module \(V\).
\end{enumerate}
Then \(\eta_A\) is a quasi-isomorphism for every finite-length module
\(A\), including nonsplit modules.
\end{lemma}

\begin{proof}
Induct on the composition length.  Choose a short exact sequence
\(0\to A'\to A\to V\to0\) with \(V\) simple.  Naturality gives a
commuting diagram of short exact sequences of complexes.  Passing to
cohomology gives a morphism between the two long exact sequences.  The
outer vertical maps are isomorphisms by induction and by the simple case;
the five lemma gives the result for \(A\).  Equivalently, the mapping
cones form a short exact sequence and the outer cones are acyclic.
\end{proof}

\begin{warning}
It is not enough to know that the three terms have the same dimensions or
that the simple-module cohomology dimensions agree.  The natural chain map
and the commuting exact diagrams are indispensable: the finite-target
induction uses connecting maps and their adjoints.
\end{warning}

\section{Quadratic phases and the corrected Fourier interface}
\label{sec:fourier}

Let \(W\) be a finite \(\F_2\)-vector space and let
\(Q:W\to\F_2\) be quadratic with nonsingular polarization
\[
 B_Q(x,y)=Q(x+y)+Q(x)+Q(y).
\]

\begin{lemma}[Affine Gauss translation]\label{lem:gauss-translate}
For every linear form \(\ell\in W^\vee\), there is a unique
\(y\in W\) with \(\ell(x)=B_Q(x,y)\) for all \(x\), and
\[
 \sum_{x\in W}(-1)^{Q(x)+\ell(x)}
 =(-1)^{Q(y)}\sum_{x\in W}(-1)^{Q(x)}.
\]
\end{lemma}

\begin{proof}
Nonsingularity of \(B_Q\) gives the unique \(y\).  Since
\(Q(x)+B_Q(x,y)=Q(x+y)+Q(y)\), translation by \(y\) gives the formula.
\end{proof}

\begin{corollary}[Gauss and zero counts]\label{cor:gauss-count}
If \(\dim W=2m\), then
\[
 \Gsum(Q):=\sum_x(-1)^{Q(x)}=\eps(Q)2^m,
 \qquad \eps(Q)\in\{+1,-1\},
\]
and
\[
 \#Q^{-1}(0)=2^{2m-1}+\eps(Q)2^{m-1}.
\]
The pair \((\dim W,\eps(Q))\) determines the isometry class of \(Q\).
\end{corollary}

\begin{definition}[Affine determinant interface]\label{def:affine-B4}
A pair of sources has the \emph{affine determinant interface} if, for
every boundary-framed lower map, every self-dual simple head, and every
linear phase produced by a scalar central pushout, the two sources have
the same affine Gauss sum
\[
 \sum_x(-1)^{Q(x)+\ell(x)}.
\]
Equivalently, one may supply:
\begin{enumerate}
\item equality of the base dimensions and base Gauss signs;
\item the common polar/edge pairing; and
\item a source-independent phase-cover theorem that converts the sign
\((-1)^{Q(y)}\) in \cref{lem:gauss-translate} into exact-image counts for
strictly smaller marked targets.
\end{enumerate}
\end{definition}

\begin{proposition}[Why the draft's B4 is insufficient]
Equality of the unshifted sums \(\Gsum(Q_1)=\Gsum(Q_2)\) does not, without
additional phase compatibility, identify all affine fibres used by the
finite-target induction.
\end{proposition}

\begin{proof}
The shifted coefficient at \(\ell=B_Q(-,y)\) differs from the unshifted
one by \((-1)^{Q(y)}\), by \cref{lem:gauss-translate}.  The unshifted
sum alone does not identify which target-side phase corresponds to a
vector with \(Q(y)=0\) or \(1\).  A marked quadratic isometry or the
phase-cover mechanism of \cref{def:affine-B4} supplies the missing data.
\end{proof}

\begin{remark}[Existing formal theorem]
The current \texttt{gq2-lean} development does not make this erroneous
inference.  Its source interface contains separate half-torsor,
separation, nondegeneracy, cocycle-cardinality, and ramified/unramified
Gauss-residue obligations, and the exact-image recursion consumes those
obligations.  The arbitrary-degree development should parameterize that
existing record rather than replace it by the four-line B1--B4 record in
the draft.
\end{remark}

\section{The local determinant quadratic form}

Let \(V\) be a nontrivial simple finite \(\F_2[G_K]\)-module whose action
factors through a finite tame group \(H\).  Let
\(q:V\to\F_2\) be an \(H\)-invariant nonsingular quadratic form, and let
\(\kappa_q^0\in H^2(V\rtimes H,\F_2)\) be an equivariant extraspecial
class restricting trivially to the zero section.  For a cocycle \(b\), put
\[
 Q^0_{K,V}([b])=\inv_K((b,\rho)^*\kappa_q^0).
\]
The existence and normalization of \(\kappa_q^0\) should be packaged as
part of the determinant datum; it is not automatic from the symbol
``\(H\)-invariant \(q\)'' unless the orthogonal action has been lifted to
the chosen extraspecial extension.

\begin{proposition}[Well-definedness and polarization]
\label{prop:local-polar}
The function \(Q^0_{K,V}\) is well-defined on \(H^1(K,V)\), and
\[
 B_{Q^0}(x,y)=\inv_K(x\smile_{b_q}y),
\]
where \(b_q\) is the polarization of \(q\).  In particular,
\(Q^0_{K,V}\) is nonsingular.
\end{proposition}

\begin{proof}
Cohomologous cocycles give graph homomorphisms
\(G_K\to V\rtimes H\) conjugate by an element of \(V\).  Inner
automorphisms act trivially on group cohomology, so their pullbacks of
\(\kappa_q^0\) agree.  This proves well-definedness.

The commutator form of the extraspecial extension is \(b_q\).  Pulling
back along the graph of \(x+y\) and comparing with the two individual
graphs gives
\[
 Q^0(x+y)+Q^0(x)+Q^0(y)=\inv_K(x\smile_{b_q}y).
\]
Since \(b_q:V\simeq V^\vee\) and local Tate duality makes the cup pairing
perfect, the polarization is nonsingular.
\end{proof}

\subsection{Unramified simple modules}

\begin{lemma}[One Hermitian line]\label{lem:hermitian-line}
Let \(E=\F_{2^{2e}}\), \(E_0=\F_{2^e}\), and
\(x^*=x^{2^e}\).  For \(a\in E_0^\times\), set
\[
 Q_a(x)=\Tr_{E_0/\F_2}(a x x^*).
\]
Then
\[
 \sum_{x\in E}(-1)^{Q_a(x)}=-2^e.
\]
\end{lemma}

\begin{proof}
The norm map \(E^\times\to E_0^\times\), \(x\mapsto xx^*\), is
surjective with fibres of size \(2^e+1\).  Hence
\[
\begin{aligned}
 \sum_{x\in E}(-1)^{Q_a(x)}
 &=1+(2^e+1)\sum_{t\in E_0^\times}
       (-1)^{\Tr_{E_0/\F_2}(at)}\\
 &=1-(2^e+1)=-2^e,
\end{aligned}
\]
because a nontrivial additive character sums to \(-1\) over
\(E_0^\times\).
\end{proof}

\begin{proposition}[Unramified local sign]\label{prop:unramified-sign}
If \(V\) is nontrivial, simple, self-dual, and unramified, then
\[
 \Gsum(Q^0_{K,V})=(-1)^n2^{n\dim(V)/2}.
\]
\end{proposition}

\begin{proof}
The finite image is cyclic of odd order.  Self-duality makes
\(E=\operatorname{End}_H(V)\) a finite field with a nontrivial adjoint
involution; write \(E=\F_{2^{2e}}\) and
\(E_0=E^{*=1}=\F_{2^e}\).  The space \(V\) is one-dimensional over
\(E\).  Local Euler--Poincar\'e and the vanishing of \(H^0\) and \(H^2\)
give
\[
 \dim_{\F_2}H^1(K,V)=n\dim_{\F_2}V,
\]
so \(H^1(K,V)\) has \(E\)-dimension \(n\).

The polar form in \cref{prop:local-polar} is the trace of a nonsingular
Hermitian form.  Diagonalize that Hermitian form over \(E/E_0\).  The
restriction of \(Q^0\) to each orthogonal \(E\)-line is invariant under
the norm-one scalars: such a scalar is an \(H\)-equivariant \(q\)-isometry,
so functoriality of the normalized extraspecial class preserves \(Q^0\).
A quadratic refinement of a nonsingular Hermitian line with this
invariance is \(Q_a\) for some \(a\in E_0^\times\).  By
\cref{lem:hermitian-line}, every line has Gauss sum \(-2^e\).  Gauss sums
multiply under orthogonal direct sum, giving
\[
 (-2^e)^n=(-1)^n2^{ne}=(-1)^n2^{n\dim(V)/2}.
\]
\end{proof}

\subsection{Ramified simple modules}

The ramified proof should be split into the following local input package.
This package is already the shape used by the existing deep-part
formalization for \(G_{\Q_2}\); all statements are field-generic.

\begin{definition}[Dyadic deep-unit package]\label{def:deep-package}
Let \(L/K\) be a finite tame splitting extension for \(V\), let
\(H=\Gal(L/K)\), and let \(e_L=v_L(2)\).  A dyadic deep-unit package
consists of:
\begin{enumerate}
\item the projective inflation--restriction isomorphism
\[
 H^1(K,V)\simeq
 \Hom_H(V^\vee,L^\times/L^{\times2});
\]
\item Hilbert orthogonality for the square-class filtration
\(\overline U_i\subset L^\times/L^{\times2}\);
\item the assertion that the ramified \(V\)-isotypic part has no middle
constituent and that
\[
 X_+:=\Hom_H(V^\vee,\overline U_{e_L+1})
\quad\text{has dimension }\frac{n\dim V}{2};
\]
\item vanishing of the normalized graph obstruction on \(X_+\), obtained
from the Shapiro--Evens orbit formula and the deep-unit Evens-norm lemma
for unramified quadratic extensions.
\end{enumerate}
\end{definition}

\begin{proposition}[Construction of the deep Lagrangian]
\label{prop:deep-lagrangian}
Assume the package of \cref{def:deep-package}.  If \(V\) is ramified,
then \(X_+\) is a totally singular Lagrangian for \(Q^0_{K,V}\).
\end{proposition}

\begin{proof}
The projective inflation--restriction isomorphism identifies \(X_+\) with
a subspace of \(H^1(K,V)\).  Hilbert orthogonality makes the cup-product
polarization vanish on \(X_+\), and the normalized Shapiro--Evens
calculation makes the quadratic function itself vanish there.  Thus
\(X_+\) is totally singular.  Its dimension is half of
\(\dim H^1(K,V)=n\dim V\), so nonsingularity of the ambient form makes it
a Lagrangian.
\end{proof}

\begin{remark}[The parity error in the draft]
Finite tame inertia in residue characteristic \(2\) always has odd order.
Thus the sentence ``when the inertia order is even'' must be deleted.
The middle-layer argument is instead: every nonzero graded unit layer has
an odd tame-inertia character, while the exceptional middle/valuation
pieces have trivial inertia; a ramified simple cannot occur in the latter.
Projectivity makes the isotypic functor exact, and Hilbert orthogonality
pairs the shallow and deep ramified constituents.
\end{remark}

\begin{proposition}[Ramified local sign]\label{prop:ramified-sign}
If \(V\) is ramified and the deep-unit package holds, then
\[
 \Gsum(Q^0_{K,V})=2^{n\dim(V)/2}.
\]
\end{proposition}

\begin{proof}
By \cref{prop:deep-lagrangian}, the nonsingular quadratic space contains a
totally singular Lagrangian.  It is therefore the split quadratic space,
whose Gauss sign is positive.  Its dimension is \(n\dim V\), giving the
formula.
\end{proof}

\begin{theorem}[Arbitrary-dyadic local Gauss theorem]
\label{thm:local-gauss}
Under the equivariant extraspecial datum and the standard local inputs of
\cref{sec:local-inputs},
\[
 \dim_{\F_2}H^1(K,V)=n\dim V
\]
and
\[
 \Gsum(Q^0_{K,V})=
 \begin{cases}
 (-1)^n2^{n\dim(V)/2},&V\text{ unramified},\\
 2^{n\dim(V)/2},&V\text{ ramified}.
 \end{cases}
\]
\end{theorem}

\begin{proof}
The dimension and nonsingularity follow from local Euler--Poincar\'e,
Tate duality, and \cref{prop:local-polar}.  Apply
\cref{prop:unramified-sign} or \cref{prop:ramified-sign} according to the
inertia action.
\end{proof}

\section{Marked Demu\v{s}kin matching without an all-degree free-relator theorem}
\label{sec:marked}

For a standard core \(P\), write
\[
 D_P=\langle\sigma,x_0,\ldots,x_n\mid P=1\rangle_{\mathrm{pro}-2}
\]
and let \(\chi_P:D_P\to\Z_2^\times\) and
\(\nu_P:D_P\to\Z_2\) be the displayed standard characters.

\begin{definition}[Marked-core certificate]\label{def:core-certificate}
A marked-core certificate for \((K,P)\) consists of:
\begin{enumerate}
\item an abstract Demu\v{s}kin isomorphism \(f:D_P\simeq D_K\) supplied by
Labute's classification;
\item a proof that the intrinsic orientation satisfies
\(\chi_K\circ f=\chi_P\);
\item an automorphism \(u\in\Aut(D_P)\) satisfying
\[
 \chi_P\circ u=\chi_P,
 \qquad \nu_K\circ f\circ u=\nu_P.
\]
\end{enumerate}
\end{definition}

\begin{proposition}[Marked matching reduction]\label{prop:marked-reduction}
To construct the automorphism \(u\) in
\cref{def:core-certificate}, it is enough to prove the following
per-standard-core statement:

\smallskip
\emph{Every automorphism of \(D_P^{\mathrm{ab}}\) preserving the canonical
orientation, the torsion relation vector, and the mod-\(2\) cup form that
occurs in the marked Smith--Witt classification lifts to an automorphism
of \(D_P\).}
\end{proposition}

\begin{proof}
Transport \(\nu_K\) through \(f\) and call the resulting character
\(\nu'\) on \(D_P\).  Local reciprocity, Smith normal form, and Witt
cancellation show that the equality of the marked data
\((C,I,\lambda,\gamma)\) produces an automorphism \(\bar u\) of
\(D_P^{\mathrm{ab}}\) preserving the orientation, relation vector, and
cup form, with \(\nu'\circ\bar u=\nu_P\).  Lift \(\bar u\) by the stated
per-core theorem.  The lift preserves \(\chi_P\) and corrects \(\nu'\), so
\(f\circ u\) is the required marked isomorphism.
\end{proof}

\begin{remark}
This reduction is strictly smaller and cleaner than the draft's proposed
``degree-\(j\) Nielsen differential plus exceptional three-dimensional
cokernel'' theorem.  The rank-three \(G_{\Q_2}(2)\) case is already proved
in the current Lean development by exactly this route: classify the
orientation-preserving action on the abelianization, lift it, and compose
with the abstract Labute isomorphism.  For the present paper, one needs the
same explicit lifting theorem for the rank-four \(M_\alpha\) and
\(N_\alpha\) cores; hyperbolic handles can then be added by standard
Nielsen moves preserving their commutator product.
\end{remark}

\section{Arithmetic branch correction for type \texorpdfstring{\(M_\alpha\)}{M-alpha}}

Let
\[
 C=\langle-1\rangle\times\langle u\rangle,
 \qquad u=(1-2^\alpha)^{-1}\in1+4\Z_2,
\]
and let \(\lambda:C\twoheadrightarrow\Z/2^r\Z\) be the marked
unramified quotient.  Write
\[
 \lambda(-1)=\epsilon2^{r-1},\qquad \lambda(u)=\eta.
\]

\begin{proposition}[The sign row is the unramified-\(i\) branch]
\label{prop:sign-excluded}
If \(\eta\) is even, then \(r=1\), \(\epsilon=1\), and \(K(i)/K\) is
unramified.  Consequently, under the standing hypothesis that
\(K(i)/K\) is ramified, \(\eta\) is necessarily odd.  The only new
\(M_\alpha\) rows are therefore the compact row \(r=0\) and the
procyclic-Frobenius row \(r\ge1\), \(\eta\in\Z_2^\times\).
\end{proposition}

\begin{proof}
If \(\eta\) is even, the image of the procyclic factor
\(\langle u\rangle\) is contained in
\(2\Z/2^r\Z\).  The element \(-1\) has order two, so its image is either
zero or \(2^{r-1}\).  Surjectivity of \(\lambda\) is possible only when
\(r=1\) and \(\lambda(-1)=1\), i.e. \(\epsilon=1\).  Then
\(\lambda\) vanishes on \(\langle u\rangle\), and
\[
 I=\ker\lambda=\langle u\rangle.
\]
The fixed field of \(\langle u\rangle\subset\Z_2^\times\) in the
\(2\)-power cyclotomic tower is obtained by adjoining \(i\): the
procyclic subgroup \(1+4\Z_2\) fixes \(\mu_4\), whereas \(-1\) acts on
\(i\) by complex conjugation.  By definition of \(I\), the quotient
\(C/I\) is the unramified cyclotomic quotient.  Hence \(K(i)/K\) is
unramified.
\end{proof}

\begin{corollary}[The \(\Q_2(\sqrt{-10})\) parameters]
\label{cor:minus10}
For \(K=\Q_2(\sqrt{-10})\), the type is \(M_2\) with
\[
 r=1,\qquad \epsilon=1,\qquad \eta=1.
\]
Thus the general procyclic word, not the sign word, is the relevant
family.
\end{corollary}

\begin{proof}
The unique unramified quadratic extension of \(K\) is
\(K(\sqrt5)\).  Since
\[
 \frac{-2}{5}=\frac{-10}{5^2}
\]
is a square in \(K\), this is also \(K(\sqrt{-2})\), a quadratic
subextension of the \(2\)-power cyclotomic tower.  The corresponding
quadratic character on
\(\Z_2^\times=\langle-1\rangle\times(1+4\Z_2)\) is nontrivial on both
\(-1\) and a generator congruent to \(5\pmod8\).  For
\(u=(-3)^{-1}\), one has \(u\equiv5\pmod8\), so
\(\lambda(-1)=\lambda(u)=1\).
\end{proof}

\section{Word certificates}

The remaining word calculations should not be represented in Lean by a
single opaque theorem for each field.  They should be checked by a generic
word evaluator and small certificates in a universal finite operator
algebra.

\begin{definition}[Word certificate]\label{def:word-certificate}
For a proposed word \(R\) and core \(P\), a word certificate consists of:
\begin{enumerate}
\item \textbf{tame specialization:} after setting every \(x_i=1\), the
word evaluates to \(1\) in \(T_{q_K}\);
\item \textbf{pro-\(2\) specialization:} after setting \(\tau=1\) and
\(\sigma_2=\sigma\), it evaluates to \(P\);
\item \textbf{Fox certificate:} for every simple tame module, explicit
invertible row and column operations reduce the evaluated two-relator
Jacobian to the prescribed normal form;
\item \textbf{Stokes certificate:} the recursively evaluated second-order
word identity is a natural chain map and is perfect on every simple
module;
\item \textbf{scalar certificate:} on \(\F_2\), the cup--Bockstein matrix
is the marked local Hilbert matrix;
\item \textbf{quadratic certificate:} on every self-dual simple module,
the extraspecial Hessian reduces to the claimed nondegenerate quadratic
normal form, together with the target-side affine phase data needed by
\cref{def:affine-B4}.
\end{enumerate}
\end{definition}

\begin{proposition}[Boundary specializations for the displayed rows]
\label{prop:word-boundary}
The tame and pro-\(2\) specialization parts of
\cref{def:word-certificate} hold for the corrected \(L\), compact
\(N_\alpha\), noncompact \(N_\alpha\), compact \(M_\alpha\), and
procyclic \(M_\alpha\) words in the uploaded draft.
\end{proposition}

\begin{proof}
Put \(\delta_i=(x_i\tau)^{\omega_2}x_i^{-1}\).

For the tame specialization, all \(x_i\) become \(1\).  Tame inertia is
pro-odd, so \(\tau^{\omega_2}=1\), whence every \(\delta_i\), every
correction word, and every hyperbolic handle becomes \(1\).  The remaining
factors are powers and commutators of powers of \(\sigma\), and the
balanced exponents cancel.

For the pro-\(2\) specialization, \(\tau=1\),
\(\sigma_2=\sigma\), and every \(\delta_i=1\).  In type \(L\), this gives
\[
 h_0=x_0^{\sigma^2}x_0,
 \qquad u_1^{-1}x_1^\sigma=x_1^{-1}x_1^\sigma=[x_1,\sigma],
\]
which is the marked \(L\)-core.  In compact type \(N\),
\[
 x_2^{-\sigma}x_2=[\sigma,x_2].
\]
In noncompact type \(N\), with \(g=x_1\sigma^{2^r}\),
\(x_2^{-g}x_2=[g,x_2]\), and the correction is \(1\).

For compact type \(M\), set \(A_0=x_0^{-1}\sigma^{-m}\).  The specialized
word is
\[
 A_0^2[A_0,x_1]\sigma^{2m}[\sigma,x_2],
\]
and \(2m=2^\alpha\), giving the marked core.  In the procyclic row,
\(R_{\mathrm{lin}}^{\pc}\) becomes the marked core.  In the auxiliary
copy,
\[
 \widehat C=\sigma^{2^r},\qquad
 \widehat A=\widehat C^{-m},
\]
and all four displayed factors commute; its total power is
\(-2m2^r+2^\alpha2^r=0\).  Hence the auxiliary copy and the final
\(D_0^2[D_0,D_1]\) specialize to \(1\).
\end{proof}

\begin{warning}[What is still absent]
The uploaded TeX says that exact Fox normal forms and ``the nonsingular
core matrices below'' have been calculated, but it does not display the
matrices, the row/column operations, or the universal second-order
identities.  Those are not expository details: they are precisely items
(3)--(6) of \cref{def:word-certificate}.  Until the expression trees and
certificates are supplied, the word theorems are not proved.
\end{warning}

\subsection{Certificate checker design}

The Lean checker should use the following reflected objects.
\begin{verbatim}
inductive PWord (Gen : Type)
  | one | gen (g : Gen) | mul | inv | conj | comm
  | zpow (a : Int) | z2pow (a : Z2) | profPow (a : ProfiniteInt)

structure FoxCertificate where
  rowOps : List ElementaryRowOp
  colOps : List ElementaryColOp
  target : FoxNormalForm
  verifies : applyOps rowOps colOps (foxJacobian word) = target.matrix

structure HessianCertificate where
  change : LinearEquiv ...
  target : QuadraticNormalForm
  verifiesPolar : ...
  verifiesQuadratic : ...
  affinePhase : PhaseCoverCertificate ...
\end{verbatim}
The evaluator should compute in finite semidirect, Heisenberg, and
extraspecial groups, using \cref{lem:omega-finite}.  The proof object is a
short list of elementary operations and polynomial/operator identities;
no external CAS result should be trusted without replay.

\section{The source-generic reconstruction theorem}

The finite-target induction should be imported from the existing
source-generic formal theorem, with its source record parameterized by
\(n\).  Its actual hypotheses are more precise than B1--B4: besides the
boundary, they include topological finite generation, scalar
\(H^1/H^2\) cardinalities, linear lift multiplicities, half-torsor
counts, exact-image recurrence, phase separation and nondegeneracy,
cocycle cardinalities, and the two Gauss-residue leaves.

\begin{theorem}[Abstract exact-image reconstruction]
\label{thm:source-abstract}
Let \(S_1,S_2\) be topologically finitely generated profinite groups with
the same marked boundary.  Suppose the two sources satisfy the same
source-data record for every boundary-framed finite target, including the
affine determinant interface of \cref{def:affine-B4}.  Then
\[
 \#\Sur(S_1,G)=\#\Sur(S_2,G)
\]
for every finite group \(G\), and \(S_1\cong S_2\).
\end{theorem}

\begin{proof}
The proof is the boundary-framed exact-image induction already formalized
for the \(G_{\Q_2}\) presentation.  Its well-founded measure is the wild
normal \(2\)-subgroup of the marked target together with the proper
exact-image strata.  The terminal scalar lane uses the boundary and
scalar cohomology leaves.  A minimal nontrivial module layer is handled by
the linear lift-torsor and adjoint obstruction data.  For a self-dual
head, the determinant pushout is converted by finite Fourier inversion
and the phase-cover identities into central covers with strictly smaller
measure.  Inclusion--exclusion closes the exact-image strata.  Summing
boundary frames gives the unframed count.  Equality of finite quotient
counts and topological finite generation give profinite reconstruction.

For arbitrary \(n\), the recursion is unchanged.  Only source-supplied
cardinalities such as
\[
 \#\Hom_{\mathrm{cont}}(S,\F_2)=2^{n+2},
 \qquad \#Z^1(S,V)=|V|^{n+1}
\]
change, and both sources supply the same values.  Thus the Lean change is
a parameterization of the leaf record, not a new induction proof.
\end{proof}

\section{Standard local inputs and formal trust boundary}
\label{sec:local-inputs}

The following inputs may initially be represented by named literature
interfaces, provided that every normalization is explicit.

\begin{longtable}{>{\raggedright\arraybackslash}p{0.24\textwidth}
                  >{\raggedright\arraybackslash}p{0.47\textwidth}
                  >{\raggedright\arraybackslash}p{0.20\textwidth}}
\toprule
Input & Exact statement used & Formal destination\\
\midrule
\endhead
Local Euler characteristic & For nontrivial simple \(V\),
\(\dim H^1(K,V)=n\dim V\), after the \(H^0,H^2\) vanishings are shown. &
\path{DyadicLocalData.eulerChar}\\
Tate duality & Perfectness and functoriality of the cup pairing, including
adjoint connecting maps. & \path{TateDualityG K 2}\\
Local reciprocity & The full \(\Z_2\)-valued geometric unramified
character and the marked cyclotomic quotient \((C,I,\lambda,\gamma)\). &
\texttt{MarkedRecip}\\
Square-class filtration & Exact graded pieces, inertia characters,
Hilbert orthogonality, and the deep-half dimension. &
\texttt{SquareClass}\\
Projectivity & A faithful ramified simple tame \(\F_2[H]\)-module is
projective, so isotypic functors and inflation--restriction are exact. &
\path{RamifiedSimple.projective}\\
Shapiro--Evens & Normalized orbit formula for the extraspecial graph
class, including involution-stabilized orbits. &
\texttt{ShapiroEvens}\\
Deep Evens norm & The degree-two Evens norm of a deep unit in an
unramified quadratic extension of dyadic fields vanishes. &
Existing deep-part theorem\\
Demu\v{s}kin classification & Abstract standard presentation and
intrinsic canonical orientation. & \texttt{LabuteClass}\\
Profinite reconstruction & Equality of finite epimorphism counts for
topologically finitely generated profinite groups implies isomorphism. &
Existing reconstruction API\\
\bottomrule
\end{longtable}

The trust boundary should not contain any assertion that a particular
explicit candidate word presents \(G_K\).  Those are the conclusions to
be proved from the reflected word certificates.

\section{Lean module graph and acceptance criteria}

A practical dependency order is:
\begin{enumerate}
\item \texttt{Dyadic/Parameters}: \(n,q_K\), semantic generators,
profinite exponents, and branch data;
\item \texttt{Dyadic/TameBoundary}: \cref{lem:tame-odd,lem:O2-tame,prop:specializations,thm:boundary};
\item \texttt{Dyadic/MarkedCore}: the reduction of \cref{sec:marked} and
per-core automorphism-lifting certificates;
\item \texttt{Dyadic/WordSyntax}, \texttt{FoxEval},
\texttt{HeisenbergEval}, and \texttt{ExtraspecialEval};
\item \texttt{Dyadic/LocalGauss}: the field-generic version of
\cref{thm:local-gauss}, reusing the current deep-part files;
\item \texttt{Dyadic/SourceDataN}: parameterized source leaves and an
adapter to the current source-generic induction;
\item one branch file for each of \(L\), compact/noncompact \(N\), and
compact/procyclic \(M\);
\item the final theorem, which invokes only
\cref{thm:certificate-main}.
\end{enumerate}

Every branch merge should satisfy:
\begin{enumerate}
\item the TeX word is generated from the same syntax tree used by Lean;
\item \texttt{\#print axioms} for the branch theorem contains only the
named local/literature interfaces;
\item no \texttt{sorryAx}, no field-specific isomorphism axiom, and no
unchecked CAS equality;
\item both boundary specializations are theorem-level equalities;
\item the Fox certificate includes explicit invertible operations;
\item the Hessian certificate proves nonsingularity and the affine phase
interface, not only the unshifted Gauss sign;
\item the \(n=1\) adapter reproduces the existing \(G_{\Q_2}\) source
record;
\item the \(\Q_2(\sqrt{-10})\) instance uses the procyclic parameters of
\cref{cor:minus10}.
\end{enumerate}

\section{Exact remaining theorem inventory}

After the corrections and reductions above, the proof is complete modulo
the following explicit leaves.

\begin{longtable}{>{\raggedright\arraybackslash}p{0.08\textwidth}
                  >{\raggedright\arraybackslash}p{0.30\textwidth}
                  >{\raggedright\arraybackslash}p{0.48\textwidth}}
\toprule
ID & Theorem/certificate & Completion criterion\\
\midrule
\endhead
MC-M & Marked automorphism lifting for the rank-four \(M_\alpha\) core &
A finite list of orientation-preserving Nielsen generators lifts the
Smith--Witt stabilizer on the abelianization; composition gives every
required marking correction.\\
MC-N & Marked automorphism lifting for the rank-four \(N_\alpha\) core &
Same criterion for the \(N\)-relation vector and cup form.\\
WC-L & Universal word certificate for \(R_L\) &
Fox, Stokes, scalar, and quadratic certificates, with the existing
\(n=1\) theorem as base and hyperbolic handle stability.\\
WC-N0 & Universal word certificate for compact \(N_\alpha\) &
In particular the unramified block must contain the invertible
\(1-S^{-1}\) term produced by \(x_2^{-\sigma}\).\\
WC-Npc & Universal word certificate for noncompact \(N_\alpha\) &
Replay the correction \(E_{r,\eta}\) symbolically and prove the claimed
cross operators for all allowed \(r,\eta\).\\
WC-M0 & Universal word certificate for compact \(M_\alpha\) &
Replay the reversed correction order and prove both projector normal
forms.\\
WC-Mpc & Universal word certificate for procyclic \(M_\alpha\) &
Replay both the linear and self-replicating copies, including all
\(T\)-dependent central terms and affine phases.\\
LG-K & Field-generic deep-unit package &
Parameterize the already formalized cochain/Evens arguments from
\(\Q_2\) to arbitrary finite dyadic \(K\), with the square-class
filtration as a named local input.\\
SD-n & Degree-\(n\) source-data adapter &
Replace hard-coded \(8\), \(|V|^2\), and the base Gauss exponent in the
current source record by the corresponding parameter fields; the
induction theorem itself remains unchanged.\\
\bottomrule
\end{longtable}

\begin{theorem}[Formal completion criterion]
When the entries in the preceding table are represented by Lean terms and
the final theorem has no axioms beyond the standard local inputs of
\cref{sec:local-inputs}, the ramified-\(i\) part of the presentation
theorem is fully formalized.  No additional mathematical induction or
finite-target classification remains.
\end{theorem}

\begin{proof}
MC-M/MC-N and the already available rank-three result give every
marked-core certificate.  The WC entries give every word certificate.
LG-K supplies \cref{thm:local-gauss}; SD-n supplies
\cref{thm:source-abstract}.  Apply \cref{thm:certificate-main} branch by
branch.  Proposition \ref{prop:sign-excluded} shows that the listed
branches exhaust the ramified-\(i\) case.
\end{proof}

\section{Recommended changes to the uploaded manuscript}

The manuscript should make the following changes before presenting the
main theorem as proved.

\begin{enumerate}
\item Replace the current B4 by \cref{def:affine-B4}, or state the exact
source-data/phase-cover theorem already used by the \(G_{\Q_2}\) proof.
\item Replace the proof architecture of the source-interface theorem by a
precise reference to the source-generic exact-image theorem and list all
fields of the source record.
\item Replace the all-degree relative Labute proof outline by
\cref{def:core-certificate,prop:marked-reduction}; prove MC-M and MC-N in
an appendix or in Lean.
\item Remove the phrase ``when the inertia order is even'' and split the
ramified local sign proof into the four components of
\cref{def:deep-package}.
\item Remove the sign-Frobenius row from the ramified-\(i\) assembly and
classify \(\Q_2(\sqrt{-10})\) by \cref{cor:minus10}.
\item Insert the actual Fox matrices, elementary operations, Hessians, and
affine phase certificates.  If they are too long for the paper, generate
an appendix from the Lean expression trees and include hashes linking the
paper formulas to those trees.
\item Until MC-M, MC-N, and the WC certificates are present, state the
explicit all-\(K\) presentation as a certificate theorem rather than an
unconditional theorem.
\end{enumerate}

\section*{Conclusion}

The revised compact relators pass the finite \(A_4\) regression that the
previous formulas failed.  The generic mathematical architecture is also
sound after replacing the four informal interfaces by the precise source
record.  The proof is not yet complete at the word level: the uploaded
draft still omits the universal Fox/Stokes/Hessian identities and the
rank-four marked automorphism-lifting theorem.  This packet isolates those
obligations, proves the surrounding reductions, removes an impossible
branch, and gives a direct route to kernel-checked completion.

\end{document}
